\section{Oscillatory integral estimates}\label{sec:oscillation}

Throughout, $\T=\R/\Z$, $\e(t)=\exp(2\pi i t)$, and Haar measure on
each torus has mass one. We use the Fourier convention
$\widehat f(k)=\int_\T f(t)\e(-kt)\,dt$, and the same convention on
$\R$. Constants denoted by $C$ are absolute and may change from line
to line. Subscripts indicate dependence on fixed data. The length
$N$ will tend to infinity only after all auxiliary data have been fixed.

We record three elementary forms of stationary phase. The first supplies
uniform control of every coefficient of a truncated quadratic wave.
The second identifies the leading coefficients of the main construction;
the third handles its small end pieces and gaps.

\begin{lemma}[Uniform quadratic formula]\label{lem:quadratic}
Let $g\in C_c^\infty(\R)$ and $\beta\in\R\setminus\{0\}$. As
$T\to\infty$, uniformly for all $u\in\R$,
\begin{equation}\label{eq:quadratic-formula}
 \int_\R g(y)\e(T\beta y^2/2-uy)\,dy
 =\frac{\e(\sgn(\beta)/8-u^2/(2T\beta))}{\sqrt{T|\beta|}}
   \left[g\left(\frac{u}{T\beta}\right)+O_{g,\beta}(T^{-1})\right].
\end{equation}
\end{lemma}
\begin{proof}
Put $y_0=u/(T\beta)$ and complete the square. Fourier inversion reduces
the bracket in \eqref{eq:quadratic-formula} to the exact expression
\[
 \int_\R \widehat g(\xi)\e(\xi y_0)
                     \e\bigl(-\xi^2/(2T\beta)\bigr)\,d\xi.
\]
For completeness, this identity follows by inserting
$\exp(-\pi r(y-y_0)^2)$, $r>0$, before applying Fourier inversion.
The inner Gaussian integral has factor
$(r-iT\beta)^{-1/2}\exp(-\pi\xi^2/(r-iT\beta))$;
its modulus is at most $(T|\beta|)^{-1/2}$. Dominated convergence
against $|\widehat g(\xi)|$ justifies $r\downarrow0$ and yields the
factor $\e(\sgn(\beta)/8)/\sqrt{T|\beta|}$.
Finally, replacing $\e(-\xi^2/(2T\beta))$ by $1$ costs at most
\[
 \frac{\pi}{T|\beta|}\int_\R \xi^2|\widehat g(\xi)|\,d\xi,
\]
independently of $y_0$.
\end{proof}

\begin{lemma}[Uniform stationary phase]\label{lem:stationary-phase}
Let $I$ be a compact interval, let $\phi$ be smooth on a neighborhood
of $I$ with $\phi''$ nowhere zero there, and let
$f\in C_c^\infty(\operatorname{int}I)$. Uniformly for $x\in[0,1]$,
\begin{equation}\label{eq:stationary-phase}
 \sqrt N\int_I f(t)\e\bigl(N(\phi(t)-xt)\bigr)\,dt
 =\frac{f(t_x)}{\sqrt{|\phi''(t_x)|}}
   \e\bigl(N(\phi(t_x)-xt_x)+\sgn(\phi'')/8\bigr)+o(1),
\end{equation}
where $t_x$ is the unique point with $\phi'(t_x)=x$ if it exists in
$I$, and the right-hand main term is zero otherwise.
\end{lemma}
\begin{proof}
Write $d=\sgn(\phi'')$. If $t_x\in I$, the change of variable
\[
 y=(t-t_x)
 \left(\frac{2d[\phi(t)-\phi(t_x)-\phi'(t_x)(t-t_x)]}
 {(t-t_x)^2}\right)^{1/2}
\]
extends smoothly across $t=t_x$ and changes the phase to
$\phi(t_x)-xt_x+dy^2/2$. Taylor's integral formula and the lower
bound on $|\phi''|$ show that this is a smooth increasing coordinate
with derivative bounded above and below, uniformly for $t_x\in I$.
All its required derivatives and those of its inverse are uniformly
bounded. The transformed amplitudes, extended by zero, have support
in a fixed compact set and uniformly bounded smooth norms, because
$\supp f$ is separated from the endpoints of $I$. At $y=0$ their
value is $f(t_x)/\sqrt{|\phi''(t_x)|}$.
Lemma~\ref{lem:quadratic}, whose error is controlled by a fixed finite
collection of smooth norms, proves the assertion uniformly in this
case. If $x\notin\phi'(I)$, the distance between $x$ and
$\phi'(\supp f)$ is bounded below uniformly; integration by parts
then gives an $O(N^{-1})$ integral. These two cases include the
endpoints of $\phi'(I)$.
\end{proof}

\begin{lemma}[Bounds with endpoint amplitudes]\label{lem:oscillatory-bounds}
Let $I$ be a compact interval, $f\in C^1(I)$, and $\phi$ be real and smooth.
If $\phi''=\Lambda\ne0$ on $I$, then for all real $k$,
\begin{equation}\label{eq:quadratic-bound}
 \left|\int_I f(t)\e(N\phi(t)-kt)\,dt\right|
 \le\frac{C}{\sqrt{N|\Lambda|}}
       \left(\norm f_\infty+\int_I|f'(t)|\,dt\right).
\end{equation}
If $N\phi'-k$ is monotone and has absolute value at least $N\rho>0$,
the same bound holds with $C/(N\rho)$ in place of
$C/\sqrt{N|\Lambda|}$.
\end{lemma}
\begin{proof}
The primitive of $\e(\pm y^2/2)$ has bounded differences on arbitrary
real intervals: inside $[-1,1]$ this is immediate, and on either
remaining half-line integration by parts gives a uniform bound.
Completing the square and rescaling therefore bounds every subinterval
integral of $\e(N\phi-kt)$ by $C/\sqrt{N|\Lambda|}$ in the quadratic
case. In the monotone case, one integration by parts bounds every such
integral by $C/(N\rho)$; monotonicity bounds the total variation of
$(N\phi'-k)^{-1}$ by $2/(N\rho)$. Integration against the amplitude
$f$, using either primitive, gives the stated estimates.
\end{proof}
